\section{Discussion}

We asked whether training a biologically constrained V1 circuit with a cell-local rule rather than
the exact gradient leaves a measurable trace in its population geometry, and which brings it closer
to mouse V1. Four conclusions follow from the confirmatory data.

\paragraph{The shared image geometry is set before learning.} The reliability-corrected agreement
with mouse V1 on 118 natural scenes was the same for untrained networks and for every learning rule,
under both objectives. In this circuit, what natural images look like to the population is
determined by the front end, the fitted cells and the measured wiring. That is consistent with
models in which V1 selectivity is largely feedforward \citep{lien2018}. It also cautions against
reading RSA agreement as evidence for any particular learning rule. An untrained, well-constrained
circuit already achieves it, so an untrained baseline and reliability correction are essential.

\paragraph{Rules differ in where they can take the circuit.} With matching in-vivo statistics as
the goal, the local rule optimized the trained statistics further than BPTT. It did so by
multiplying recurrent weights roughly 90-fold and driving interneurons hard. This solution
generalized poorly to natural images: dimensionality, decodability and sparseness moved away from
mouse values. BPTT stopped earlier and retained mouse-like dimensionality and decodability.
Crucially, at matched loss the two rules produced the same geometry. The difference is therefore
not an intrinsic signature of local credit assignment but a consequence of which regions of
weight space each rule can reach. Local rules ignore network feedback in their gradient estimates
and can keep pushing in a direction that the exact gradient would correct. This argues for adding
mechanisms that bound such drift, such as weight or rate homeostasis, when local rules are used to
fit data, and for comparing learning rules at matched performance.

\paragraph{On a behavioral task the representation barely moves.} Training change detection on
eight images changed natural-scene geometry little under any rule. The local rule solved the task
better than BPTT at the shared learning rates, and with symmetric feedback it raised
dimensionality towards the mouse value. These effects are small compared with the
mouse--model gap.

\paragraph{The spike-history term.} Unlike in our earlier abstract tasks, removing $\beta$ from the
eligibility trace did not change any geometry measure here. Its importance appears to depend on
whether the task requires credit assignment through the cells' adaptation dynamics.

\paragraph{Relation to the specification's hypotheses.} We did not find that local learning yields
representations ``significantly more'' similar to mouse V1 than BPTT. On the preregistered endpoint
the rules were indistinguishable, and on several secondary measures BPTT was closer to the mouse
under data matching.

\paragraph{Limitations.}
\begin{itemize}
\item \emph{Front end and fixed inputs.} The LGN model and the LGN$\to$V1 weights are simplified
and fixed. Plastic thalamocortical weights would give learning more leverage over geometry.
\item \emph{Rates, not spikes; 20-ms bins.} Fano factors and spike-timing statistics cannot be
compared with the recordings.
\item \emph{Scale and cells.} The circuit has 1{,}000 cells sampled with replacement from about 300
fits, and represents one L2/3-like column without laminar structure.
\item \emph{Learning rates.} Rates were shared across rules and chosen on a small development grid.
BPTT was unstable at the chosen data-matching rate late in training. The change-detection
recurrent rate was not tuned. Matched-performance comparisons address but do not remove this
dependence.
\item \emph{Mouse data.} We used six sessions (627 V1 units) recorded during passive viewing. Units
were pooled across animals into a pseudo-population, and putative cell classes were assigned by
waveform only.
\item \emph{Multiplicity.} The secondary contrasts were consistent across all 10 networks but do not
survive Holm correction over the full secondary family.
\item \emph{Disattenuation.} When a model's RDM reliability is low (0.14 for local data-matching
networks), the corrected RSA is noisy.
\end{itemize}

